\cotefile{161-1}
\scanpage{4}

% loc: upper right
\margin{Foncteurs adjoints\\[2pt] [demander Schanuel]}

% loc: upper left
Soit

% loc: upper centre
\diagram{\begin{tikzcd}[ampersand replacement=\&,column sep=large]
E\arrow[r,shift left=.7ex,"u"]\& F\arrow[l,shift left=.7ex,"v"]\&[-.5em]
\begin{aligned}i&:\operatorname{id}_{E}\longrightarrow vu\\ j&:uv\longrightarrow\operatorname{id}_{F}\end{aligned}
\end{tikzcd}}

% loc: below diagram
un couple de foncteurs adjoints.

% loc: a, opening
a) Soit $E'$ une sous-catégorie \deleted{\illegible{1}} pleine de $E$, \deleted{\unsure{et stable par}} Notons que pour $X,Y\in\operatorname{Ob}E'$, l'application

% loc: a, Hom display
\[
\begin{aligned}
\operatorname{Hom}_{E'}(X,Y)&\longrightarrow\operatorname{Hom}_{F}(uX,uY)\quad(\simeq\operatorname{Hom}(X,vuY))\\[-2pt]
\rotatebox{90}{$\simeq$}&\\[-2pt]
\operatorname{Hom}_{E}(X,Y)&
\end{aligned}
\]

% loc: a, explanation
est déduite de l'application $i(Y):\deleted{\unsure{vu}}Y\longrightarrow vu(Y)$ comme $\operatorname{Hom}(X,i(Y))$. Donc \deleted{\unsure{dire}}\inserted{pour} que \deleted{\unsure{ceci}}\inserted{ce soit injectif} \deleted{\unsure{signifie}}\inserted{(resp.)} pour $\forall X,Y\in\operatorname{Ob}E'$ \deleted{\unsure{signifie}}\inserted{i.e. $u|E'$ \deleted{pl.} fidèle (resp. pl. fidèle)} il suffit que $\forall Y\in\operatorname{Ob}E'$ \deleted{\unsure{on ait}}, $i(Y):\deleted{\unsure{soit}}Y\to vu(Y)$ soit \deleted{\unsure{un isom.}}\inserted{un mono} (resp. un \uline{\unsure{isom|mono}}). C'est aussi nécessaire si on sait que $vu(E')\subset\overline{E'}$ (image ess. de l'inclusion $E'\subset\unsure{\overline E|E}$\unsure{)}.

% loc: a, transition
\deleted{Donc Donc \unsure{l'on a}} On trouve donc

% loc: a, condition 1
\[
u|E'\ \text{pl. fidèle et }vu(E')\subset\overline{E'}\quad\Longleftrightarrow\quad
\forall\,Y\in\operatorname{Ob}E',\quad i(Y):Y\xrightarrow{\sim}vu(Y)\tag*{\textcircled{1}}\]

% loc: under condition 1
\margin{donner un exemple montrant que cette condition n'est pas surabondante, avec $E'$ réduit à un élément}

% loc: b, opening
b) Soit $F'$ une sous-catégorie pleine de $F$. Notons que pour $P,Q\in\operatorname{Ob}F'$, l'app.

% loc: b, Hom display
\[
\begin{aligned}
\operatorname{Hom}_{F'}(P,Q)&\longrightarrow\operatorname{Hom}(vP,vQ)\quad(\simeq\operatorname{Hom}(uvP,Q))\\[-2pt]
\rotatebox{90}{$\simeq$}&\\[-2pt]\operatorname{Hom}_{F}(P,Q)&
\end{aligned}
\]

% loc: b, explanation
est déduite de $j(P):uv(P)\longrightarrow P$ par $\operatorname{Hom}(j(P),Q)$. Donc on trouve pour \inserted{pour que} $v|F'$ soit fidèle \inserted{(pl. fidèle)}, il suffit que $\forall P\in\operatorname{Ob}F'$, $j(P):\inserted{uv(P)\to P}$ soit un épi. (resp. un isom). Donc on trouve

% loc: b, condition 2
\[
v|F'\ \text{pl. fidèle et }uv(F')\subset\overline{F'}\quad\Longleftrightarrow\quad
\forall P\in\operatorname{Ob}F',\ j(P):uv(P)\xrightarrow{\sim}P.\tag*{\textcircled{2}}\]

% loc: below condition 2, barred two-line draft
\deleted{Si $\overline{F'}$ est stable par quotients, on a}\\
\deleted{$v|F'$ fidèle et $uv(F')\subset\overline{\unsure{F|F'}}\Longleftrightarrow\forall P\in\operatorname{Ob}F',\ j(P):uv(P)\to P$ un}

% loc: c, correspondence
c) L'ensemble des $E'$ \inserted{strictement pleines $\subset E$} satisfaisant (1), et des $F'$ \inserted{strictement pleines} satisfaisant (2), sont en correspondance biunivoque, en posant

% loc: c, images
\[
\begin{aligned}
\alpha(E')&=\text{image essentielle de }u|E'=\overline{u(E')}\\
\beta(F')&=\text{image essentielle de }v|F'=\overline{v(F')}.
\end{aligned}
\]

% loc: c, equivalences
Si $E'$ et $F'$ \deleted{\unsure{se correspondent ainsi}, alors \unsure{$u|E'$ et $v'$}}\inserted{sont des sous-catégories \inserted{str.} pleines \deleted{\unsure{l'une l'autre}} de $E$ resp. $F$}, les conditions suivantes sont équivalentes :

% loc: c, lines 1 to 5
1) $E'$ satisfait \textcircled{1}, $F'$ satisfait \textcircled{2}, et $E'$ et $F'$ se correspondent comme dessus

2) $u(E')\subset F'$, $v(F')\subset E'$, et $u'$ \deleted{est pl. fidèle}\inserted{et $v'$} sont pl. fidèles

3) \rule{1.1em}{.4pt}\quad\rule{2.1em}{.4pt}\quad et $u'$ une équivalence

4) \rule{3.8em}{.4pt}\quad et $v'$ une équivalence

5) \rule{3.8em}{.4pt}\quad et $u'$ et $v'$ pl. fid.

% loc: c, right of last lines
$\unsure{6|5})\ \ldots\ldots$, et $\left\{\begin{aligned}
\forall\,\deleted{\unsure{X}}Y\in\operatorname{Ob}E',\ i(Y)&:Y\xrightarrow{\sim}vuY\\
\forall\,P\in\operatorname{Ob}F',\ j(P)&:uvP\xrightarrow{\sim}P
\end{aligned}\right.$

% archive sheet number: 3
